\documentclass[../../../include/open-logic-section]{subfiles}

\begin{document}

\olfileid{sth}{ordinals}{basic}
\olsection{الخصائص الأساسية للأعداد الترتيبية}

لاحظنا أن الأعداد الترتيبية الأولى هي الأعداد الطبيعية. والسبب الرئيس
لتطوير نظرية للأعداد الترتيبية هو توسيع مبدأ الاستقراء الذي يصح على
الأعداد الطبيعية. وسنبلغ هذا عبر سلسلة من النتائج الأولية.

\begin{lem}\ollabel{ordmemberord}
كل !!{element} في عدد ترتيبي هو عدد ترتيبي.
\end{lem}

\begin{proof}
ليكن~$\alpha$ عددًا ترتيبيًا و$b \in \alpha$. لما كان~$\alpha$
متعديًا، كان~$b \subseteq \alpha$. ومن ثم ترتب~$\in$ المجموعة~$b$
ترتيبًا جيدًا، لأن~$\in$ ترتب~$\alpha$ ترتيبًا جيدًا.

ولرؤية أن~$b$ متعد، افترض~$x \in c \in b$. عندئذ~$c \in \alpha$
لأن~$b \subseteq \alpha$. ومرة أخرى، لما كان~$\alpha$ متعديًا، كان
$c \subseteq \alpha$، ومن ثم~$x \in \alpha$. ولذلك
$x, c, b \in \alpha$. لكن~$\in$ ترتب~$\alpha$ ترتيبًا جيدًا، فتكون
$\in$ علاقة متعدية على~$\alpha$ بموجب
\olref[sth][ordinals][wo]{wo:strictorder}. ولما كان~$x \in c \in b$،
كان~$x \in b$. وبالتعميم، $c \subseteq b$.
\end{proof}

\begin{cor}\ollabel{ordissetofsmallerord}
$\alpha = \Setabs{\beta \in \alpha}{\beta \text{ عدد ترتيبي}}$، لكل
عدد ترتيبي~$\alpha$.
\end{cor}

\begin{proof}
ينتج مباشرة من~\olref{ordmemberord}.
\end{proof}

الخلاصة التقريبية للنتيجتين الرئيسيتين التاليتين، وهما
\olref{ordinductionschema} و\olref{ordtrichotomy}، أن الأعداد
الترتيبية نفسها مرتبة ترتيبًا جيدًا بعلاقة الانتماء:

\begin{thm}[الاستقراء المتناهي التجاوز]\ollabel{ordinductionschema}
لكل صيغة~$\phi(x)$:	
\[
	\text{إذا كان }\exists \alpha \phi(\alpha)\text{، فإن }
	\exists \alpha(\phi(\alpha)
	\land (\forall \beta \in \alpha) \lnot \phi(\beta))
\]
حيث تقيد المكممات المعروضة ضمنًا بالأعداد الترتيبية.
\end{thm}

\begin{proof}
%\olref{ordinductionschema1}. 
افترض~$\phi(\alpha)$ لعدد ترتيبي ما~$\alpha$. إذا كان
$(\forall \beta \in \alpha) \lnot \phi(\beta)$، فقد انتهينا. وإلا
فلتكن
$X = \Setabs{\beta \in \alpha}{\phi(\beta)}$؛ عندئذ
$X \neq \emptyset$. ولما كان~$\in$ يرتب~$\alpha$ ترتيبًا جيدًا، كان
لـ$X$ !!{element} أصغر بحسب~$\in$، وليكن~$\gamma$. وبحسب تعريف~$X$،
يكون~$\phi(\gamma)$، كما أن~$\gamma$ عدد ترتيبي بموجب
\olref{ordmemberord}. وإذا كان~$\beta \in \gamma$، فإن
$\beta \in \alpha$ لأن~$\alpha$ متعد، ولا يمكن أن يكون~$\beta \in X$
بسبب أصغرية~$\gamma$؛ ومن ثم~$\lnot\phi(\beta)$. وبالتعميم،
$(\forall \beta \in \gamma)\lnot\phi(\beta)$.
%\olref{ordinductionschema2}. Suppose there is some ordinal
%	$\gamma$ such that $\lnot\phi(\gamma)$. Then by
%	\olref{ordinductionschema1}, there is an $\in$-minimal ordinal
%	$\alpha$ for which $\lnot\phi(\alpha)$. So $(\forall \beta <
%	\alpha) \phi(\beta)$, rendering the antecedent of the conditional
%	false.
\end{proof}
\noindent
لاحظ أنه يمكننا بالقدر نفسه التعبير عن~\olref{ordinductionschema}
بالمخطط:
\[
\text{إذا كان }\forall \alpha((\forall \beta \in \alpha)\phi(\beta) \lif 
\phi(\alpha))\text{، فإن }\forall \alpha\phi(\alpha)
\]
بمجرد أخذ~$\lnot\phi(\alpha)$ في~\olref{ordinductionschema}، ثم إجراء
معالجات منطقية أولية.

\begin{thm}[التثليث]\ollabel{ordtrichotomy} 
$\alpha \in \beta \lor \alpha = \beta \lor \beta \in \alpha$،
لأي عددين ترتيبيين~$\alpha$ و$\beta$.
\end{thm}

\begin{proof}
البرهان باستقراء مزدوج، أي باستعمال~\olref{ordinductionschema} مرتين.
لنقل إن~$x$ \emph{قابل للمقارنة} مع~$y$ إذا وفقط إذا كان
$x \in y \lor x = y \lor y \in x$.

لأجل الاستقراء، افترض أن كل عدد ترتيبي في~$\alpha$ قابل للمقارنة مع
\emph{كل} عدد ترتيبي. ولأجل استقراء إضافي، افترض أن~$\alpha$ قابل
للمقارنة مع كل عدد ترتيبي في~$\beta$. سنبيّن أن~$\alpha$ قابل للمقارنة
مع~$\beta$. وبالاستقراء على~$\beta$، ينتج أن~$\alpha$ قابل للمقارنة مع
كل عدد ترتيبي؛ ثم بالاستقراء على~$\alpha$ يكون \emph{كل} عدد ترتيبي
قابلًا للمقارنة مع \emph{كل} عدد ترتيبي، كما هو مطلوب. ويكفي أن نفترض
$\alpha \notin \beta$ و$\beta \notin \alpha$، وأن نبيّن
$\alpha = \beta$.

لإثبات~$\alpha \subseteq \beta$، ثبّت~$\gamma \in \alpha$؛ وهذا عدد
ترتيبي بموجب~\olref{ordmemberord}. ولذلك يكون~$\gamma$، بموجب فرض
الاستقراء الأول، قابلًا للمقارنة مع~$\beta$. لكن إذا كان
$\gamma = \beta$ أو كان~$\beta \in \gamma$، لنتج~$\beta \in \alpha$
(مع استعمال كون~$\alpha$ متعديًا إذا لزم)، خلافًا لافتراضنا؛ ومن ثم
$\gamma \in \beta$. وبالتعميم، $\alpha \subseteq \beta$.

ويبيّن استدلال مماثل تمامًا، باستعمال فرض الاستقراء الثاني، أن
$\beta \subseteq \alpha$. ولذلك~$\alpha = \beta$.
\end{proof}\noindent
وبناء على ذلك سنكتب أحيانًا~$\alpha <\beta$ بدلًا من
$\alpha \in \beta$، لأن~$\in$ تسلك مسلك علاقة ترتيب. ولا أسباب
عميقة لذلك تتجاوز الألفة، وكون كتابة~$\alpha \leq \beta$ أسهل من
كتابة~$\alpha \in \beta \lor \alpha = \beta$.
\footnote{كان يمكننا أن نكتب~$\alpha \mathrel{\underline{\in}} \beta$؛
لكن ذلك غير معياري إطلاقًا.}

وفيما يأتي نتيجتان سريعتان لنتائجنا الأخيرة، تستعمل أولاهما ترميزنا
الجديد:

\begin{cor}\ollabel{ordordered}
إذا كان~$\exists \alpha\phi(\alpha)$، فإن
$\exists \alpha(\phi(\alpha) \land
\forall \beta(\phi(\beta) \lif \alpha \leq \beta))$. وفوق ذلك، لكل
أعداد ترتيبية~$\alpha, \beta, \gamma$، يصح كل من
$\alpha \notin \alpha$ و
$\alpha \in \beta \in \gamma \lif \alpha \in \gamma$.
\end{cor}

\begin{proof}
مثل~\olref[wo]{wo:strictorder} تمامًا.
\end{proof}

\begin{prob}
أكمل «الاستدلال المماثل تمامًا» في برهان
\olref[sth][ordinals][basic]{ordtrichotomy}.
\end{prob}

\begin{cor}\ollabel{corordtransitiveord}
تكون~$A$ عددًا ترتيبيًا إذا وفقط إذا كانت~$A$ مجموعة متعدية من الأعداد
الترتيبية.
\end{cor}

\begin{proof}
من \emph{اليسار إلى اليمين}: بموجب~\olref{ordmemberord}. ومن
\emph{اليمين إلى اليسار}: إذا كانت~$A$ مجموعة متعدية من الأعداد
الترتيبية، فإن~$\in$ ترتب~$A$ ترتيبًا جيدًا بموجب
\olref{ordinductionschema} و\olref{ordtrichotomy}.
\end{proof}

قلنا، على سبيل التفسير، إن~\olref{ordinductionschema} و
\olref{ordtrichotomy} تخبراننا بأن~$\in$ ترتب الأعداد الترتيبية ترتيبًا
جيدًا. غير أنه ينبغي أن نكون \emph{حذرين جدًا} من مثل هذا الادعاء،
بسبب النتيجة الآتية:

\begin{thm}[مفارقة بورالي--فورتي]\ollabel{buraliforti}
لا توجد مجموعة لجميع الأعداد الترتيبية.
\end{thm}

\begin{proof}
برهانًا بالخلف، افترض أن~$O$ مجموعة جميع الأعداد الترتيبية. إذا كان
$\alpha \in \beta \in O$، كان~$\alpha$ عددًا ترتيبيًا بموجب
\olref{ordmemberord}، ومن ثم~$\alpha \in O$. ولذلك تكون~$O$ متعدية،
فتكون عددًا ترتيبيًا بموجب~\olref{corordtransitiveord}. ومن ثم
$O \in O$، وهذا يناقض~\olref{ordordered}.
\end{proof}
\noindent
سميت هذه النتيجة باسم~\citeauthor{Burali-Forti1897}. لكن كانتور هو
أول من رأى بوضوح، في رسالة إلى ددكند عام~1899، \emph{التناقض} في
افتراض وجود مجموعة لجميع الأعداد الترتيبية. وكما يشرح فان هاينورت:
\begin{quote}
  رأى بورالي--فورتي نفسه أن التناقض يثبت، بطريقة \emph{البرهان بالخلف}،
  أن الترتيب الطبيعي للأعداد الترتيبية ليس إلا ترتيبًا جزئيًا.
  \citep[p.~105]{Heijenoort1967}
\end{quote}
إذا نحينا خطأ بورالي--فورتي جانبًا، استطعنا تلخيص ما سبق هكذا: الأعداد
الترتيبية مجموعات يرتب الانتماء كلًا منها منفردًا ترتيبًا جيدًا، ويرتبها
مجتمعة ترتيبًا جيدًا (من دون أن تؤلف مجتمعةً مجموعة).

وختامًا، فيما يأتي مزيد من الخصائص الأساسية للأعداد الترتيبية الناتجة
من~\olref{ordinductionschema} و\olref{ordtrichotomy}.

\begin{prop}
كل متتالية تنازلية تمامًا من الأعداد الترتيبية متناهية.
\end{prop}

\begin{proof}
كل متتالية لا نهائية تنازلية تمامًا من الأعداد الترتيبية
$\alpha_0 > \alpha_1 > \alpha_2 > \ldots$ لا تملك عنصرًا أصغريًا
بحسب~$<$، خلافًا لـ\olref{ordinductionschema}.
\end{proof}

\begin{prop}\ollabel{ordinalsaresubsets}
$\alpha \subseteq \beta \lor \beta \subseteq \alpha$، لأي عددين
ترتيبيين~$\alpha, \beta$.
\end{prop}

\begin{proof}
إذا كان~$\alpha \in \beta$، كان~$\alpha \subseteq \beta$ لأن~$\beta$
متعد. وبالمثل، إذا كان~$\beta \in \alpha$، كان
$\beta \subseteq \alpha$. وإذا كان~$\alpha = \beta$، كان كل من
$\alpha \subseteq \beta$ و$\beta \subseteq \alpha$. ولذلك انتهينا
بموجب~\olref{ordtrichotomy}.
\end{proof}

\begin{prop}\ollabel{ordisoidentity}
$\alpha = \beta$ إذا وفقط إذا كان~$\ordeq{\alpha}{\beta}$، لأي عددين
ترتيبيين~$\alpha, \beta$.
\end{prop}

\begin{proof}
الأعداد الترتيبية ترتيبات جيدة؛ ولذلك تنتج النتيجة مباشرة من التثليث
(\olref[basic]{ordtrichotomy}) ومن
\olref[iso]{wellordnotinitial}.
\end{proof}

\begin{prob}\ollabel{probunionordinalsordinal}
أثبت أنه إذا كان كل عنصر من~$X$ عددًا ترتيبيًا، فإن~$\bigcup X$ عدد
ترتيبي.
\end{prob}

\end{document}
